\section{Target Generator}

\subsection{Scopo}

Il \textit{Target Generator} è un componente utilizzato durante il training per
trasformare le bounding box del ground truth in target compatibili con la
formulazione FCOS.

La Detection Head produce infatti una predizione per ogni posizione spaziale
delle feature map della FPN, ma durante l'addestramento è necessario stabilire
che cosa ogni posizione dovrebbe predire. Il Target Generator esegue proprio
questa associazione.

La pipeline può essere vista come:

\[
\text{Ground Truth}
\rightarrow
\text{Target Generator}
\rightarrow
\text{Target}
\rightarrow
\text{Loss}
\]

Il Target Generator non fa parte del modello neurale e viene utilizzato
solamente durante il training. Durante l'inference il ground truth non è
disponibile e quindi questo componente non viene utilizzato.

\subsection{Target prodotti}

Per ogni livello della FPN, il Target Generator produce tre target:

\begin{itemize}
    \item \texttt{positive}: indica se una location della feature map
    corrisponde a un oggetto;
    
    \item \texttt{ltrb}: contiene le quattro distanze dalla location ai lati
    della bounding box;
    
    \item \texttt{centerness}: indica quanto la location si trova vicino al
    centro della bounding box.
\end{itemize}

In questo modo, per ogni location positiva, il modello deve imparare a
predire:

\[
(l,t,r,b)
\]

e il relativo valore di \textit{centerness}.

\subsection{Conversione delle location}

La prima operazione consiste nel trasformare una posizione
$(x,y)$ della feature map nelle corrispondenti coordinate dell'immagine
originale.

Per una feature map con stride $s$, il centro della location viene calcolato
come:

\[
x_{\text{img}}=(x+0.5)s,
\qquad
y_{\text{img}}=(y+0.5)s.
\]

Il termine $0.5$ serve a considerare il centro della cella della feature map.

Questa operazione è necessaria perché le bounding box del ground truth sono
espresse nelle coordinate dell'immagine originale, mentre le predizioni della
Detection Head sono associate alle location delle feature map.

\subsection{Verifica della presenza della location nel box}

Per ogni location viene verificato se il punto trasformato nelle coordinate
dell'immagine cade all'interno di una bounding box.

Se:

\[
x_1 \leq x \leq x_2
\]

e

\[
y_1 \leq y \leq y_2,
\]

la location è considerata appartenente al bounding box e può quindi diventare
un candidato positivo.

Se la location si trova fuori dal box, quel ground truth non può essere
assegnato alla location.

\subsection{Calcolo delle distanze LTRB}

Quando una location si trova all'interno di una bounding box, vengono
calcolate le distanze tra la location e i quattro lati del box:

\[
l=x-x_1
\]

\[
t=y-y_1
\]

\[
r=x_2-x
\]

\[
b=y_2-y.
\]

Il vettore risultante è:

\[
(l,t,r,b).
\]

Questa è la rappresentazione utilizzata da FCOS per la regressione della
bounding box: invece di predire direttamente $(x_1,y_1,x_2,y_2)$, il modello
predice la distanza della location dai quattro lati del box.\footnote{Zhi Tian
et al., \textit{FCOS: Fully Convolutional One-Stage Object Detection},
ICCV 2019.}

\subsection{Assegnazione in base alla scala}

Non tutte le feature map sono adatte a rappresentare oggetti della stessa
dimensione. Per questo il Target Generator utilizza un diverso
\textit{regression range} per ogni livello della FPN.

Nel progetto:

\[
\begin{array}{c|c}
\text{Stride} & \text{Regression range} \\
\hline
8   & [0,64) \\
16  & [64,128) \\
32  & [128,256) \\
64  & [256,512) \\
128 & [512,\infty)
\end{array}
\]

Per ogni location viene considerata la massima delle quattro distanze:

\[
m=\max(l,t,r,b).
\]

Il bounding box può essere assegnato a un determinato livello solo se $m$
rientra nel range associato a quel livello.

Questo meccanismo permette di distribuire gli oggetti di dimensioni diverse
tra i diversi livelli della FPN e riduce l'ambiguità tra le feature map.
L'uso di un diverso intervallo di regressione per ogni livello è una parte
fondamentale della formulazione FCOS.\footnote{Zhi Tian et al.,
\textit{FCOS: Fully Convolutional One-Stage Object Detection}, ICCV 2019.}

\subsection{Gestione di più bounding box compatibili}

Una location può trovarsi contemporaneamente all'interno di più bounding box.
In questo caso può esserci ambiguità su quale oggetto debba essere utilizzato
come target.

Il progetto risolve questa situazione scegliendo il bounding box con area
minore tra quelli compatibili.

\[
B^*=\arg\min_{B_i} \operatorname{Area}(B_i)
\]

Questa strategia segue la regola utilizzata nella formulazione FCOS per
risolvere le location appartenenti a più ground-truth box.\footnote{Zhi Tian
et al., \textit{FCOS: Fully Convolutional One-Stage Object Detection},
ICCV 2019.}

\subsection{Calcolo della Center-ness}

Una volta scelto il bounding box, viene calcolato il target di
\textit{centerness} a partire dalle quattro distanze LTRB:

\[
\text{centerness}
=
\sqrt{
\frac{\min(l,r)}{\max(l,r)}
\cdot
\frac{\min(t,b)}{\max(t,b)}
}.
\]

Una location vicina al centro del bounding box avrà un valore vicino a $1$,
mentre una location vicina ai bordi avrà un valore più basso.

Il motivo è che le location lontane dal centro tendono a produrre bounding box
di qualità inferiore. FCOS introduce quindi la center-ness come una predizione
separata utilizzata per ridurre il peso di queste detection durante
l'inference.\footnote{Zhi Tian et al., \textit{FCOS: Fully Convolutional
One-Stage Object Detection}, ICCV 2019.}

\subsection{Costruzione finale dei target}

Dopo aver eseguito tutte le verifiche, ogni location della feature map viene
classificata come positiva o negativa.

Per una location negativa:

\[
\texttt{positive}=0
\]

e i target di regressione e center-ness vengono impostati a zero.

Per una location positiva:

\[
\texttt{positive}=1
\]

e vengono memorizzati:

\[
\texttt{ltrb}=(l,t,r,b)
\]

e

\[
\texttt{centerness}\in[0,1].
\]

Il risultato finale è quindi una struttura di target della stessa
organizzazione spaziale della feature map, che verrà successivamente
confrontata con le predizioni della Detection Head attraverso la funzione di
loss.